\section{Tesla 路线：视觉、BEV、Transformer 和自研芯片}

前面说明早期路线的问题，本章进入 Tesla。郎咸朋认为，2018 年前后自动驾驶路线发生关键转向：Tesla 坚持不用高精地图和激光雷达，转向纯视觉、BEV、Transformer 和自研芯片。它的核心不是某一个名词，而是把多个相机的信息放到同一个空间表征中，让系统统一理解车辆周围世界。

\begin{figure}[H]
\centering
\includegraphics[width=0.96\textwidth]{tesla-updimension.png}
\caption{Tesla 的升维解法：用车队数据、视觉、芯片和规模化训练解决系统问题。自制概念图，依据 00:04:30--00:25:06 对谈内容整理。}
\end{figure}

\begin{knowledgebox}{读图：升维不是堆更多补丁}
图中 Fleet Data、Vision、Chip、Training 和 FSD 连成一条线。Tesla 的“升维”在于不再修补每个后融合错误，而是改变表示方式和系统边界：用统一空间、统一模型、车队数据和自研算力解决问题。
\end{knowledgebox}

\subsection{为什么用视觉}

本节解释视觉路线。Elon Musk 常用“人只有眼睛也能开车”来解释纯视觉，这适合大众传播；技术层面更重要的是，图像信息密度高，包含颜色、纹理、语义、边界和隐含几何。通过多相机视差、连续帧运动和模型学习，系统可以从图像中推断三维空间、速度和趋势。

\begin{figure}[H]
\centering
\includegraphics[width=0.92\textwidth]{camera-lidar-tradeoff.png}
\caption{Camera 与 LiDAR 的取舍：LiDAR 直接给距离，Camera 给语义和低成本规模。自制概念图，依据 00:25:06--00:28:46 对谈内容整理。}
\end{figure}

\begin{knowledgebox}{读图：直接测距不等于信息更多}
LiDAR 的优点是点上距离准；Camera 的优点是像素密度和语义信息丰富。郎咸朋的判断是，激光雷达只给 128 行左右的点云，而摄像头给数百万像素的彩色世界。自动驾驶需要的不只是距离，还包括对象、语义、运动意图和环境上下文。
\end{knowledgebox}

\subsection{BEV 和 Transformer 解决什么}

本节解释 BEV。BEV 即 Bird's Eye View，鸟瞰视角。朴素做法是每个摄像头先各自识别目标，再把结果后融合；问题是不同相机可能漏检、重叠、尺度不一致，后融合容易出错。BEV 的思想是先提取多相机特征，把信息统一到一个空间表征里，再一次性识别人、车、车道和障碍物。Transformer 在这里是通用建模架构，关键是它服务于统一空间理解。

\begin{importantbox}{BEV 的关键不是鸟瞰图好看，而是避免后融合错误}
后融合是“各相机先独立判断，再拼结果”；BEV 更接近“先汇集多相机特征，再在统一空间里判断”。这能减少不同来源结果互相打架的问题，让对象一致性更好。
\end{importantbox}

\subsection{前融合与后融合：为什么思想比算子更重要}

本节把郎咸朋对 BEV 的解释再展开。早期多摄像头方案常见的朴素做法，是每个相机各自跑一套前向感知：前摄像头识别前方车道和车辆，侧摄像头识别侧向目标，后摄像头识别后方目标；随后工程系统再把这些检测结果融合到同一个坐标系。问题在于，每个相机的结果都可能错、漏、重复或尺度不一致，后处理需要判断哪个相机可信、两个半辆车是否同一辆车、遮挡边缘如何拼合。这会把大量精力耗在“消除后融合误差”上。

BEV 的升维之处，是把融合提前到特征层：先把多相机图像提取成特征，再在统一空间中推理对象、车道和可行驶区域。这样做不是单纯“把图拼成大图”，而是让模型在同一个空间里同时使用多个视角的信息。老师强调，Transformer 只是后来服务这个思想的网络结构，真正关键的是先看到了自动驾驶行业的本质问题：多视角感知必须统一建模，而不是每个视角局部判断后再拼。

\begin{knowledgebox}{课堂提示：BEV 的 lesson 可以迁移到很多 AI 系统}
当一个系统有多个信息源时，后融合常常带来一致性问题。更本质的解法，是在合适的表征空间里做早期融合，让模型一次性利用所有上下文。这和多模态模型、Agent 工具状态、机器人 VLA 都有相似思路。
\end{knowledgebox}

\begin{figure}[H]
\centering
\includegraphics[width=0.92\textwidth]{av-terms-map.png}
\caption{自动驾驶关键术语地图：BEV、Transformer、E2E、VLM、World Model 和 RL 的关系。自制概念图，依据 00:20:06--01:26:40 对谈内容整理。}
\end{figure}

\begin{knowledgebox}{读图：术语之间不是并列清单}
BEV 和 Transformer 主要解决空间表征与建模；E2E 把输入到轨迹的映射交给模型；VLM 增加语义理解；World Model 预测未来状态；RL 让系统在反馈中优化。读这张图时，要把它看成一条从感知到决策再到闭环学习的链条。
\end{knowledgebox}

\subsection{为什么自研芯片}

本节看芯片。多相机 BEV 和大模型计算需要大量车端算力。Tesla 早期从 NVIDIA 芯片切到自研 FSD 芯片，逻辑是：算法、传感器和芯片共同设计，才能以较低成本获得有效算力。访谈里提到 Tesla 传感器+芯片成本约 1000 美元级别，重点不是精确数字，而是说明量产路线必须把算力、成本和算法适配一起设计。

\begin{knowledgebox}{术语消化：TOPS、ASIC、有效算力}
\noindent\begin{tabularx}{\textwidth}{p{0.18\textwidth}p{0.36\textwidth}X}
\toprule
术语 & 含义 & 本期中的作用 \\
\midrule
TOPS & 每秒万亿次运算，常用于车载 AI 芯片算力指标 & 数值高不等于实际模型跑得好。 \\
ASIC & 专用集成电路，为特定算法/负载定制 & Tesla 可用自研算法和芯片互相适配。 \\
通用芯片 & 可支持多家公司和多种算法的处理器 & 灵活，但未必为某一算法最优。 \\
有效算力 & 实际被模型利用的算力 & 比标称 TOPS 更接近工程结果。 \\
\bottomrule
\end{tabularx}
\end{knowledgebox}

\subsection{本章小结}

Tesla 路线的本质，是把自动驾驶从昂贵传感器和地图确定性，转向视觉、统一空间表征、车队数据、自研芯片和大规模训练。这是自动驾驶第一次明确走向“模型能力系统”。

